
\documentclass[a4paper,11pt]{article}
\usepackage{fontspec}
\setmainfont{Noto Sans}
\usepackage[french]{babel}
\usepackage{microtype}
\usepackage[top=9mm,bottom=8mm,left=11mm,right=11mm]{geometry}
\usepackage{xcolor}
\usepackage{tikz}
\usetikzlibrary{arrows.meta,calc,positioning}
\usepackage[most]{tcolorbox}
\usepackage{ulem}
\pagestyle{empty}

\definecolor{Green}{HTML}{319B58}
\definecolor{GreenDark}{HTML}{17613A}
\definecolor{GreenLight}{HTML}{EAF6EC}
\definecolor{GreenLine}{HTML}{8FC9A3}
\definecolor{GreenBorder}{HTML}{B9D3C3}
\definecolor{Orange}{HTML}{D8842F}
\definecolor{OrangeDark}{HTML}{9A5716}
\definecolor{OrangeLight}{HTML}{FFF1DF}
\definecolor{OrangeLine}{HTML}{E7B77E}
\definecolor{OrangeBorder}{HTML}{E7C7A6}
\definecolor{Violet}{HTML}{7A5AA6}
\definecolor{VioletDark}{HTML}{563A7A}
\definecolor{VioletLight}{HTML}{F2ECF8}
\definecolor{VioletLine}{HTML}{B9A5D0}
\definecolor{VioletBorder}{HTML}{D3C4E2}
\definecolor{Ink}{HTML}{172A32}
\definecolor{Grey}{HTML}{64757D}
\definecolor{Red}{HTML}{D52F35}
\definecolor{RedLight}{HTML}{FCEBED}

\tcbset{
 enhanced,boxrule=.6pt,arc=2.8mm,
 left=3.7mm,right=3.7mm,top=1.75mm,bottom=1.75mm,
 boxsep=1mm,coltext=Ink
}
\newtcolorbox{greenbox}[1][]{colback=GreenLight,colframe=GreenBorder,#1}
\newtcolorbox{orangebox}[1][]{colback=OrangeLight,colframe=OrangeBorder,#1}
\newtcolorbox{violetbox}[1][]{colback=VioletLight,colframe=VioletBorder,#1}
\newtcolorbox{warnbox}[1][]{colback=RedLight,colframe=Red,
 left=2.7mm,right=2.7mm,top=1.25mm,bottom=1.25mm,#1}

\begin{document}
\begin{tikzpicture}[remember picture,overlay]
\draw[draw=VioletBorder,line width=.9pt,rounded corners=1.7mm]
 ([xshift=1mm,yshift=-1mm]current page.north west)
 -- ([xshift=-1mm,yshift=-1mm]current page.north east)
 -- ([xshift=-1mm,yshift=1mm]current page.south east)
 -- ([xshift=1mm,yshift=1mm]current page.south west)--cycle;
\fill[Violet] ([xshift=1mm,yshift=-1mm]current page.north west)
 rectangle ([xshift=5.2mm,yshift=1mm]current page.south west);
\end{tikzpicture}

\vspace*{1mm}\hspace*{6mm}
{\color{VioletDark}\fontsize{25}{27}\selectfont\bfseries CONDITIONNEL PASSÉ}

\vspace{0.5mm}\hspace*{6mm}
{\color{Grey}\large Le conditionnel passé}

\vspace{2.5mm}\hspace*{6mm}
\begin{minipage}{0.91\textwidth}

\begin{violetbox}
\textbf{\color{VioletDark}CONSTRUCTION}\\[1mm]
\centering
{\Large \textbf{AUXILIAIRE AU CONDITIONNEL PRÉSENT + PARTICIPE PASSÉ}}\\[1mm]
\small
\textit{j'aurais parlé \quad je serais parti(e)}
\end{violetbox}

\vspace{2.5mm}
\begin{center}
\begin{tikzpicture}[>=Latex]
\draw[Ink!65,line width=.8pt] (0,0)--(15.1,0);
\node[below=2mm] at (2.0,0) {\small PASSÉ};
\node[below=2mm] at (7.5,0) {\small PRÉSENT};
\node[below=2mm] at (13.5,0) {\small FUTUR};
\draw[Violet,line width=2.5pt,->] (2.6,0)--(6.4,0);
\draw[VioletDark,line width=1.1pt,->] (7.5,1.0)--(7.5,.18);
\node[above=3.8mm] at (7.5,1.0) {\bfseries\color{VioletDark}MAINTENANT};
\node[above=3mm,align=center,text width=6.2cm] at (4.5,0)
{\small\bfseries action passée non réalisée, hypothétique ou dépendante d'une condition};
\end{tikzpicture}
\end{center}

\vspace{1.5mm}
\begin{violetbox}
\textbf{\color{VioletDark}À QUOI SERT-IL ?}\\[1mm]
\begin{itemize}\setlength{\itemsep}{0.4mm}
\item exprimer une \textbf{action passée qui ne s'est pas réalisée} ;
\item exprimer un \textbf{regret} ou un reproche ;
\item présenter une hypothèse portant sur un fait passé ;
\item rapporter une information passée avec \textbf{prudence} ou réserve.
\end{itemize}
\end{violetbox}

\vspace{2mm}
\begin{violetbox}
\textbf{\color{VioletDark}EXEMPLES}\\[1mm]
\textbf{Si j'avais su, je serais venu plus tôt.}\\
\textbf{J'aurais aimé vous prévenir.}\\
\textbf{Elle aurait pu nous aider.}\\
\textbf{Ils seraient partis avant notre arrivée.}
\end{violetbox}

\vspace{2mm}
\begin{minipage}{0.485\textwidth}
\begin{violetbox}[height=3.25cm]
\textbf{\color{VioletDark}AVEC « SI »}\\[1mm]
Pour une hypothèse irréelle dans le passé :\\[1mm]
\centering
\textbf{SI + plus-que-parfait}\\
$\downarrow$\\
\textbf{conditionnel passé}\\[1mm]
\raggedright
\textit{Si j'avais su, je serais venu.}
\end{violetbox}
\end{minipage}\hfill
\begin{minipage}{0.485\textwidth}
\begin{warnbox}[height=3.25cm]
\textbf{\color{Red}ATTENTION}\\[1mm]
Après \textbf{si}, on n'emploie normalement pas le conditionnel dans cette construction.\\[1mm]
\textbf{Si j'avais su...}\\
\textcolor{Grey}{et non : \sout{Si j'aurais su...}}
\end{warnbox}
\end{minipage}

\vspace{2mm}
\begin{violetbox}
\textbf{\color{VioletDark}À RETENIR}\\[1mm]
Le conditionnel passé présente une action passée comme \textbf{hypothétique, non réalisée,
regrettée ou dépendante d'une condition}. Avec \textbf{être}, le participe passé s'accorde avec le sujet.
\end{violetbox}

\vspace{1.5mm}\centering
{\small\color{Grey}\textbf{Repères fréquents :} si, sans, sinon, heureusement, malheureusement, à ta place, j'aurais dû, j'aurais aimé...}

\end{minipage}
\end{document}
